
\begin{abstract}
Recent advances in deep monocular visual \ac{SLAM} have achieved impressive accuracy and dense reconstruction capabilities, yet their robustness to scale inconsistency in large-scale indoor environments remains largely unexplored. Existing benchmarks are limited to room-scale or structurally simple settings, leaving critical issues of intra-session scale drift and inter-session scale ambiguity insufficiently addressed. To fill this gap, we introduce the \textit{ScaleMaster Dataset}, the first benchmark explicitly designed to evaluate scale consistency under challenging scenarios such as multi-floor structures, long trajectories, repetitive views, and low-texture regions. We systematically analyze the vulnerability of state-of-the-art deep monocular visual SLAM systems to scale inconsistency, providing both quantitative and qualitative evaluations. Crucially, our analysis extends beyond traditional trajectory metrics to include a direct map-to-map quality assessment using metrics like Chamfer distance against high-fidelity 3D ground truth. Our results reveal that while recent deep monocular visual \ac{SLAM} systems demonstrate strong performance on existing benchmarks, they suffer from severe scale-related failures in realistic, large-scale indoor environments. By releasing the ScaleMaster dataset and baseline results, we aim to establish a foundation for future research toward developing scale-consistent and reliable visual \ac{SLAM} systems.

\end{abstract}


